\begin{claim}{A2.3}{Proved}{Identification of the elliptic logarithm and the verified zero formula}
For the Weierstrass equation
\[
y^2=x^3-3t^3(t-1)x+2t^4(t-1)^2,
\]
the section $[\mu]$ is a generator of its rational section group and
equals one of $P_*=(0,\sqrt2t^2(t-1))$ and $-P_*$.
In particular $\wp(\mu,\tau)=0$ on every regular fibre.
Along the branch in \eqref{eq:A-s-branch}, with
$x_0=1-1/t$, the exact branch is
\begin{equation}\label{eq:A-DI-branch}
\mu=\frac{1-\tau}{2}
 +\frac{i\sqrt6}{3\pi}
 \frac{x_0^{1/4}\,{}_3F_2(1/3,2/3,1;3/4,5/4;x_0)}
 {{}_2F_1(1/12,5/12;1;1-x_0)}.
\end{equation}
Consequently $\mu-(2-\tau)/3$ has exact order two at the
order-three point in $r^3=t$.
\end{claim}
\begin{proof}
We include the comparison, since a matching $j$-map alone would not
identify a section.  Write $F$ for the elliptic scale, not the
threefold projection.  It is specified by
\[
F^2=-\frac{3t(t-1)E_4(\tau)}{\pi^2E_6(\tau)},\qquad
F/t\longrightarrow i\sqrt3/\pi
\]
at the distinguished cusp.  The local orders are two at $r=0$ and
one at $s=0$; hence the square root exists on the simply connected
marked cover.  Its transformation laws are
$F(g_1z)=-F(z)/\tau(z)$ and $F(g_2z)=F(z)/\tau(z)$.
These follow first after squaring by modularity, and the signs are
fixed by those local orders.  Direct substitution in the cubic
identifies its regular fibres with $\C/(\Z+\tau\Z)$ by
\[
x=F^2\wp(\nu,\tau),\qquad y=\tfrac12F^3\wp'(\nu,\tau).
\]
The transformation laws of $\mu$ make this a single-valued section
$P$.  Its coordinates are meromorphic at the two finite special
points and at the cusp, so are rational functions on $\Pone$.

Every rational section $Q$ admits a global holomorphic elliptic
logarithm $\nu$ on the marked cover.  Indeed the exponential map
from the trivial complex line bundle to this smooth elliptic family
is a covering map with locally uniformly discrete lattice; pullback
over a simply connected base lifts the section.  The apparent poles
in its projective Weierstrass coordinates at a finite point are
meromorphic poles of finite order: the coordinates of $Q$ are
rational in $t$, while $F$ is holomorphic with the finite orders
just computed.  Multiplying the three homogeneous coordinates by
a common power of the local uniformizer and cancelling their
common zero gives holomorphic coordinates not all zero at the
centre.  Their limiting point lies on the closed projective
Weierstrass family.  This is the valuative criterion for properness
over the local analytic discrete valuation ring; it is not an
assertion that an arbitrary punctured-disc holomorphic map to a
compact space has a removable singularity.
We can choose this logarithm bounded at the distinguished cusp.
To verify this last assertion, put $v=1/t$, $X=v^2x$, $Y=v^3y$.
The equation becomes
$Y^2=X^3-3(1-v)X+2(1-v)^2$.  A section cannot meet its nodal point
$(X,Y)=(1,0)$ at $v=0$: substituting $X=1+O(v),Y=O(v)$ leaves
nonzero coefficient one in the defining equation at order $v$.
On the smooth nodal locus the coordinate
$u=\exp(2\pi i\nu)$ is a nonzero finite coordinate, as follows
directly from
\[
X=1+\frac{12u}{(u-1)^2},\qquad
Y=\frac{12\sqrt3u(u+1)}{(u-1)^3}.
\]
The normally convergent Weierstrass Fourier series extends this
coordinate to a nearby smoothing.  The inverse function theorem
therefore writes the section as a holomorphic unit $u(v)$, also
near $u=1$, the point at infinity.  A logarithm of that unit fixes
the bounded lift.  Any two such lifts differ by an integer, since
an additional nonzero multiple of $\tau$ is unbounded.

For this lift, write the integral transformation errors as
\[
\nu(g_1z)=-\nu/\tau+\alpha_1+\beta_1(\tau-1)/\tau,
\quad
\nu(g_2z)=\nu/\tau+\alpha_2-\beta_2/\tau.
\]
Bounded cusp regularity and $g_1g_2=g_0^{-1}$ imply
$\alpha_2=-\beta_1$, $\beta_2=\alpha_1+\beta_1$.
Putting $n=-\alpha_1-2\beta_1$, $a=\alpha_1+\beta_1$, we find that
$\nu-(n\mu+a)$ satisfies the homogeneous law of AC26 Lemma 3.10
and is cusp-regular.  We recall a direct verification of the degree
used here.  The function
$F_h=E_4(\tau)^2E_6(\tau)^{1/2}/\Delta(\tau)$ on the marked cover
has orders two and one at $z_1,z_2$, and has a simple pole at the
cusp.  The square root exists because its zeros on this cover have
even order.  Modularity and these local orders fix the two signs:
$F_h(g_1z)=-F_h/\tau$, $F_h(g_2z)=F_h/\tau$.
For any holomorphic solution of the same homogeneous law, the
quotient by $F_h$ is an invariant meromorphic function of $t$.
A possible pole at the order-three point would have order at most
$2/3$ in $t$, and one at the order-four point at most $1/4$; its
integral order is therefore nonnegative.  At the cusp regularity
requires this quotient to vanish.  Thus the solution sheaf is
exactly $\OO_B(-p_0)$, with no nonzero global sections.
We obtain $\nu=n\mu+a$ and hence $Q=nP$.  The two distinct
finite values of $\mu$ exclude $nP=0$ for $n\ne0$ by cusp boundedness;
there is no torsion.  This proves that $P$ generates the section group.

We identify $P_*$ with a primitive generator by an intersection
calculation.  The rational elliptic surface and its fibre types
$\mathrm{IV}^*,\mathrm{III},I_1$ are those in Engel Section 2.
At $t=1$, the blowup chart $x=(t-1)X,y=(t-1)Y$ gives
$Y^2=(t-1)X^3-3t^3X+2t^4$; $P_*$ meets its exceptional
multiplicity-one component at $(X,Y)=(0,\sqrt2)$, away from the
identity component.  At $t=0$, two successive charts
$x=t^2X,y=t^2Y$ give
$Y^2=t^2X^3-3t(t-1)X+2(t-1)^2$; $P_*$ meets the smooth exceptional
component $Y=-\sqrt2$ at $X=0$.  These are the two nonidentity
simple components of the $E_6$ configuration.  The section is
disjoint from $O$ in every chart, including infinity.

Here the height calculation can be made directly, without assuming
primitivity.  Orthogonally project $Q-O$ in
$\operatorname{NS}(S)\otimes\Q$ off $O$, the fibre $F_S$, and the
nonidentity components.  The inverse intersection matrices are
$-C_{E_6}^{-1}$ and $-1/2$.  Their diagonal entries for the two
nonidentity simple components are $-4/3$ and $-1/2$.
Here we may number the nonidentity components so that
\[
C_{E_6}=\begin{pmatrix}
2&-1&0&0&0&0\\
-1&2&-1&0&-1&0\\
0&-1&2&-1&0&0\\
0&0&-1&2&0&0\\
0&-1&0&0&2&-1\\
0&0&0&0&-1&2
\end{pmatrix},\qquad
(C_{E_6}^{-1})_{44}=(C_{E_6}^{-1})_{66}=4/3.
\]
Since $K_S=-F_S$, adjunction gives $Q^2=O^2=-1$.  Expanding the
square of this projection gives
\[
h(Q)=2+2(Q\cdot O)-c_0(Q)-c_1(Q),\quad
c_0(Q)\in\{0,4/3\},\ c_1(Q)\in\{0,1/2\}.
\]
The projection is additive: the chord-and-tangent law makes
$Q+R-(Q+R)_{\rm section}-O$ linearly equivalent to a vertical
divisor, which the projection kills.  Its negative square therefore
satisfies $h(nQ)=n^2h(Q)$.  A section meets a single simple
component, since its base parameter has order one; consequently
the listed local corrections exhaust the possibilities.
For $Q\ne O$ we have $Q\cdot O\ge0$, so $h(Q)\ge1/6$.
For $P_*$ the explicit incidences give $h(P_*)=1/6$.
Writing $P_*=nP$ now gives $1/6=n^2h(P)\ge n^2/6$, and hence
$n=\pm1$.  Inverting the displayed $E_6$ Cartan matrix and the
chord-and-tangent formulas is included in the exact computation log.
This proves $\wp(\mu,\tau)=0$.

We checked the original theorem on page 399 of
\citet{EichlerZagier1982}, as well as equation (6) of
\citet{DukeImamoglu2008}.  With the normalized $\Delta$, their
coefficient is indeed $144\pi i\sqrt6$, their constant is
$\log(5+2\sqrt6)/(2\pi i)$, and the integrand is
$(\sigma-\tau)\Delta(\sigma)/E_6(\sigma)^{3/2}$ on the vertical
path with its continued branch.  Thus the formula in Lead 3 is
correct with its indicated sign and period ambiguity.  Theorem 1
of Duke--Imamoglu gives \eqref{eq:A-DI-branch}: we take one minus
their $z_0$, rather than $z_0-\tau$, because the latter is unbounded
at the distinguished cusp.  The chosen branch approaches
$(1-i)/2$ at $\tau=i$ and has lower imaginary cusp value.
These conditions identify the original bounded lift.  In particular
its linear coefficient is
$i\sqrt6/(3\pi\,{}_2F_1(1/12,5/12;1;1))$, agreeing with
\eqref{eq:A-mu-value}.

Finally exact elliptic addition gives
\[
x(2P_*)=9t^2/8,\quad
y(2P_*)=\frac{t^2(32-5t)}{16\sqrt2},\quad
x(3P_*)=\frac{16}{81}(t-1)(5t-32).
\]
The last expression has nonzero limit $512/81$ at $t=0$.
If $\nu=\mu-(2-\tau)/3$, then $3\nu$ is the local logarithm of
$3P$ modulo its exact period.  Since $F$ has order two and
$\wp(3\nu)=(3\nu)^{-2}+O(\nu^2)$, the equality
$x(3P)=F^2\wp(3\nu)$ forces $\nu$ to have exact order two.
It also gives its leading coefficient
$\nu_2=\pm3F_2/(16\sqrt2)$, where $F_2=\lim_{r\to0}F/r^2$;
the sign is the explicitly stated elliptic-logarithm branch.
\end{proof}
